\chapterbackmatter{Solutions to Weinberg Exercises}
\ExerciseEditorialNote

\begin{enumerate}
  \WeinbergSolution{1}
  For a constant pseudoscalar background, the quadratic fermion operator
  contains a scalar mass \(m\) and a pseudoscalar mass \(g\phi\).  In the
  Hermitian pseudoscalar convention the latter is written
  \(ig\phi\gamma _5\); the factor \(i\) may instead be absorbed into the
  definition of the coupling used in the question.  A constant chiral
  rotation combines the two masses into
  \[
    \rho^2(\phi)=m^2+g^2\phi^2 .
  \]
  Equivalently, multiplication of the Dirac operator by its conjugate
  converts its determinant to a determinant of
  \(-\partial^2+\rho^2\).  There is no gauge field here, so this constant
  change of variables has no anomalous term.

  The scalar fluctuation determinant is independent of the background,
  since the scalar theory is otherwise free, and can be absorbed into the
  vacuum-energy counterterm.  The fermion determinant has the form of the
  last term of Eq.~\eqref{eq:16.2.16}, with \(M^2(\phi)\) replaced by
  \(\rho^2(\phi)\).  Thus, in the normalization used in this chapter,
  \[
    \boxed{
    V_{\rm eff}(\phi)
      =\lambda_R+\frac12 M_R^2\phi^2+\frac{h_R}{24}\phi^4
       -\frac{\rho^4(\phi)}{32\pi^2}
        \ln\frac{\rho^2(\phi)}{\kappa^2}}
  \]
  through one loop, up to a scheme-dependent finite polynomial in
  \(\phi^2\).  Here
  \(\rho^2(\phi)=m_R^2+g_R^2\phi^2\) to the required accuracy.
  The scale \(\kappa\) merely makes the logarithm dimensionless.

  Although no \(\phi^4\) interaction was specified in the classical
  Lagrangian, the determinant contains divergences proportional to
  \(1\), \(\phi^2\), and \(\phi^4\).  Renormalization therefore requires
  the vacuum, scalar-mass, and quartic counterterms displayed above; one
  may impose \(h_R=0\) at a chosen subtraction scale, but it does not stay
  zero at every scale.  If the coupling in the question is interpreted
  literally without the Hermiticity factor \(i\), the same calculation
  gives the analytic continuation \(g^2\mapsto-g^2\).

  \WeinbergSolution{2}
  Use collective indices \(i=(n,x)\), \(j=(m,y)\), and so on, with
  repeated indices denoting both summation and spacetime integration.
  Write
  \[
    \Delta_{ij}=W_{,ij},\qquad
    \Gamma_{ij}=\Gamma_{,ij}.
  \]
  Differentiating the Legendre relations gives
  \[
    \Gamma_{ia}\Delta_{aj}=-\delta_{ij}.
  \]
  Differentiation with respect to \(J_k\), followed by multiplication by
  \(\Delta\), gives
  \[
    \boxed{W_{,ijk}
      =\Delta_{ia}\Delta_{jb}\Delta_{kc}\Gamma_{abc}.}
  \]
  The sign is positive because the full connected two-point function in
  this chapter is already
  \(\Delta=-\Gamma^{(2)-1}\).

  Differentiating once more, and using the three possible choices of the
  differentiated external propagator, yields
  \[
    \boxed{\begin{aligned}
    W_{,ijkl}
      ={}&\Delta_{ia}\Delta_{jb}\Delta_{kc}\Delta_{ld}
      \bigl[
        \Gamma_{abcd}
        +\Gamma_{abe}\Delta_{ef}\Gamma_{fcd}\\
        &\hspace{38mm}
        +\Gamma_{ace}\Delta_{ef}\Gamma_{fbd}
        +\Gamma_{ade}\Delta_{ef}\Gamma_{fbc}
      \bigr].
    \end{aligned}}
  \]
  The first formula is the tree consisting of one full 1PI three-point
  vertex and three full external propagators.  At four points the first
  term is one full 1PI four-point vertex.  The other three terms are the
  \(s\)-, \(t\)-, and \(u\)-type trees made from two full 1PI
  three-point vertices joined by one full propagator; four full external
  propagators are attached in every term.  These are precisely the
  connected trees generated by \(\Gamma\).

  \WeinbergSolution{3}
  In six dimensions the renormalizable classical potential must be
  allowed to contain the vacuum, tadpole, and mass terms as well as the
  cubic interaction:
  \[
    V_0(\phi)=\lambda+t\phi+\frac12m^2\phi^2+\frac{g}{6}\phi^3 .
  \]
  After the shift by a constant background, the quadratic fluctuation
  mass is
  \[
    \mu^2(\phi)=m^2+g\phi .
  \]
  The one-loop scalar contribution in Euclidean form is
  \[
    \mathcal I_6(x)
      =\frac12\int\frac{d^6k_E}{(2\pi)^6}\ln(k_E^2+x),
      \qquad x=\mu^2(\phi).
  \]
  Four derivatives remove the ultraviolet divergence:
  \[
    \mathcal I_6^{(4)}(x)
      =-3\int\frac{d^6k_E}{(2\pi)^6}
        \frac1{(k_E^2+x)^4}
      =-\frac{1}{2(4\pi)^3x}.
  \]
  Integrating four times gives
  \[
    \mathcal I_6(x)
      =-\frac{x^3}{12(4\pi)^3}\ln\frac{x}{\kappa^2}
        +A+Bx+Cx^2+Dx^3.
  \]
  Since \(x=m^2+g\phi\), the polynomial is exactly absorbed by the four
  counterterms already displayed.  Hence
  \[
    \boxed{\begin{aligned}
    V_{\rm eff}(\phi)={}&
      \lambda_R+t_R\phi+\frac12m_R^2\phi^2+\frac{g_R}{6}\phi^3\\
      &-\frac{[m_R^2+g_R\phi]^3}{12(4\pi)^3}
        \ln\frac{m_R^2+g_R\phi}{\kappa^2}
    \end{aligned}}
  \]
  to one-loop order, again up to a finite polynomial fixed by the
  subtraction prescription.  When \(\mu^2(\phi)<0\), continuation through
  the logarithm produces an imaginary part.  It signals that the
  homogeneous background is unstable, not that the exact effective
  potential has become complex.

  \WeinbergSolution{4}
  Regard the fields and currents as columns.  The source term is
  \(\int d^4x\,\phi^{\mathsf T}J\).  If
  \(\phi'=M\phi\), it is left unchanged by the contragredient
  transformation
  \[
    \boxed{J'=M^{-\mathsf T}J},
    \qquad
    J'_n=\sum_m(M^{-1})_{mn}J_m .
  \]
  Changing integration variables in the functional integral, and assuming
  that both the action and the measure are invariant, therefore gives
  \[
    Z[M^{-\mathsf T}J]=Z[J],
    \qquad
    W[M^{-\mathsf T}J]=W[J].
  \]
  Differentiation shows that the classical field transforms covariantly,
  \[
    \phi_{M^{-\mathsf T}J}=M\phi_J.
  \]
  On inverting this relation one obtains
  \(J_{M\phi}=M^{-\mathsf T}J_\phi\).  Substitution in the Legendre
  transform now gives
  \[
    \begin{split}
    \Gamma[M\phi]
      &= -\int d^4x\,(M\phi)^{\mathsf T}
          M^{-\mathsf T}J_\phi
         +W[M^{-\mathsf T}J_\phi]\\
      &= -\int d^4x\,\phi^{\mathsf T}J_\phi+W[J_\phi]
       =\Gamma[\phi].
    \end{split}
  \]
  Thus every finite linear symmetry of the action and measure is inherited
  by the quantum effective action.
\end{enumerate}
